\section{第40届复赛第6题：交变磁场中导体球涡电流悬浮与热效应}
\subsection{题干}
半径 $b$ 的圆线圈置于水平面，电流 $I(t)=I_0\cos\omega t$。重量 $G$、半径 $a$、电导率 $\sigma$ 的导体球球心在线圈轴线上方 $h$ 处，且 $\delta\ll a\ll b$。求理想导体球磁矩和面电流、球心磁场、悬浮所需 $I_0$ 与平均焦耳热功率。
\begin{figure}[htbp]\centering\begin{tikzpicture}[>=Stealth,scale=.88]
\draw[thick] (-2.2,0) ellipse (2.2 and .45);\draw[thick] (-2.2,0) arc[start angle=180,end angle=360,x radius=2.2,y radius=.45];
\draw[->] (0,0) -- (0,3.1) node[above] {$z$};\shade[ball color=gray!30] (0,2.15) circle(.42);\node[right] at (.45,2.15) {$a,\ G$};
\draw[<->] (.55,.12) -- (.55,1.73) node[midway,right] {$h$};\node[below] at (0,-.65) {载流线圈半径 $b$，$I(t)=I_0\cos\omega t$};
\end{tikzpicture}
\caption{第40届复赛第6题原题图6a的TikZ重绘。}\end{figure}
\subsection{参考解答整理}
\paragraph{（1）理想导体球的磁矩与面电流。}
理想导体球内总磁场为零，故
\[\boxed{\boldsymbol m=-\frac{2\pi a^3}{\mu_0}\boldsymbol B_e,\qquad \boldsymbol i=-\frac{3B_e}{2\mu_0}\sin\theta\,\hat{\boldsymbol\phi}.}\]
\paragraph{（2）线圈在球心处产生的磁场。}
线圈轴线上球心处
\[B_0=\frac{\mu_0I_0b^2}{2(b^2+h^2)^{3/2}}.\]
\paragraph{（3）导体球悬浮所需的电流振幅。}
周期平均安培力平衡给出
\[\boxed{I_0=\sqrt{\frac{4G(b^2+h^2)^4}{3\pi a^3b^4\mu_0h}}}.\]
\paragraph{（4）涡电流的周期平均热功率。}
用趋肤层均匀体电流近似并消去 $\delta$，平均热功率为
\[\boxed{\overline P=\frac{3\pi a^2b^4I_0^2}{8(b^2+h^2)^3}\sqrt{\frac{2\omega\mu_0}{\sigma}}}.\]
\subsection{转写依据}
题干依据 \texttt{CPHO40-Q06\_transcription\_v1.md}，答案依据第16至19页与 \texttt{CPHO40-Q06\_formula\_check\_v1.csv}。图6a按原题图重绘。本题四问分值为12、4、11、13分。
